\subsection{Base lógica para la inclusión y exclusión}

Decidimos qué datos entran con los tres criterios que pide el curso: relevancia para el objetivo del proyecto (predecir el BCS), calidad de los datos y restricciones técnicas.

\begin{table}[H]
  \centering
  \small
  \caption{Fuentes que entran y que no entran a la preparación.}
  \label{tab:seleccion}
  \begin{tabularx}{\textwidth}{>{\raggedright\arraybackslash}p{0.34\textwidth}l>{\raggedright\arraybackslash}X}
    \toprule
    Fuente & ¿Entra? & Razón \\
    \midrule
    Fotos de la visita al CAETEC del 14/10/2026 & Sí & Son la entrada del modelo de visión. Son 355 fotos de 69 vacas sin contar copias. \\
    \addlinespace
    Bitácora: columnas ID, hora y BCS & Sí & El BCS de los expertos es la etiqueta. \\
    \addlinespace
    Bitácora: columnas Foto y Carpeta & No & Son ligas de Drive. La relación entre foto y vaca ya viene en las carpetas. \\
    \addlinespace
    Conjunto de datos de VMS 3 (\nolinkurl{vms-3.csv}) & No & Solo 13 de las 69 vacas con BCS están en este conjunto y no es entrada del modelo de imágenes. \\
    \addlinespace
    Columna ``Evento BCS'' de VMS 3 & No & No coincide con el BCS de los expertos y todavía no sabemos qué mide (véase 2.4). \\
    \addlinespace
    CSV de ordeño & No & Solo 12 vacas con BCS tienen datos de ordeño y son de marzo a julio de 2025, más de un año antes de las fotos. \\
    \bottomrule
  \end{tabularx}
\end{table}

Otras decisiones de esta tarea:

\begin{itemize}
  \item No hicimos pruebas de significancia ni de correlación entre campos, porque la entrada del modelo son imágenes y no una tabla de atributos.
  \item No hicimos muestreo ni pesos de atributos. Tenemos pocas vacas y las fotos únicas pesan cerca de 0.79 GB, así que usamos todas las fotos disponibles.
  \item No recolectamos datos adicionales en esta fase.
  \item Si el socio formador responde qué mide ``Evento BCS'', volvemos a evaluar si se usa.
\end{itemize}

\subsubsection*{Cómo obtuvimos las cifras}

Todas las cifras de esta fase salen de correr completo el notebook \nolinkurl{01_dp_bcs_problem.ipynb} sobre las fotos originales (477 archivos), la bitácora (\nolinkurl{etiquetas_bsc.xlsx}) y el archivo \nolinkurl{vistas.csv} con las vistas marcadas a mano. La corrida terminó sin errores y el índice que generó coincide con el índice de referencia que ya estaba en el bucket (333 filas, mismas columnas y ninguna diferencia). Es decir, el notebook reproduce exactamente el dataset del bucket. El notebook busca la carpeta \nolinkurl{data/} subiendo desde su ubicación, así que no depende de una ruta fija.
